\input{estilo}
\begin{document}
\begin{figura}\begin{tikzpicture}[y=0.95cm,
    reloj/.style={font=\small\bfseries, text=segundo, inner sep=1.5pt},
    marcado/.style={reloj, circle, draw=segundo, line width=0.8pt, inner sep=1pt}]
  \foreach \x/\p in {0/P, 4/Q, 8/R} {
    \draw (\x,0) -- (\x,7.6);
    \node[font=\small, above] at (\x,7.6) {proceso \var{\p}};
  }
  \foreach \l/\i/\x/\y/\lado in {p/1/0/0.7/left, p/2/0/2.7/left, p/3/0/4.8/left, p/4/0/6.8/left,
      q/1/4/0.75/right, q/2/4/1.8/right, q/3/4/2.75/right, q/4/4/3.8/below right,
      q/5/4/4.8/right, q/6/4/5.8/right, q/7/4/6.75/right,
      r/1/8/0.7/right, r/2/8/2.75/right, r/3/8/4.7/right, r/4/8/6.75/right} {
    \node[evento] (\l\i) at (\x,\y) {};
    \node[etiq, \lado=2pt] at (\l\i) {\var{\l}\textsubscript{\i}};
  }
  \draw[onda] (p1) -- (q2);
  \draw[onda] (q1) -- (p2);
  \draw[onda] (q1) -- (r4);
  \draw[onda] (q4) -- (r3);
  \draw[onda] (q5) -- (p4);
  \draw[onda] (r2) -- (q7);
  % los valores del reloj que anotó el profesor
  \foreach \e/\v/\pos/\dist in {p1/1/below left/4pt, q1/1/below left/3pt, q2/2/above left/3pt, q3/3/left/4pt, q5/5/below left/3pt, q6/6/left/4pt,
      r1/1/right/13pt, r2/2/right/13pt}
    \node[reloj, \pos=\dist] at (\e) {\v};
  \node[marcado, left=5pt] at (q4) {4};
  \node[marcado, right=15pt] at (r3) {5};
  \node[reloj] at (7.2,4.95) {4};
  \node[reloj, below=6pt] at (8,0) {0};
\end{tikzpicture}\end{figura}
\end{document}
